\documentclass[11pt]{article}
\usepackage{fullpage,graphicx,hyperref,amsmath,algorithm,algpseudocodex}

\newcommand{\fig}[2]{\begin{figure}[h]\begin{center}%
  \includegraphics[scale=#2]{#1.pdf}\end{center}%
  \end{figure}}
\newcommand{\ceil}[1]{\left\lceil #1\right\rceil}
\newcommand{\floor}[1]{\left\lfloor #1\right\rfloor}


\begin{document}

\begin{center}\Large\bf 
CS 473 (Spring 2025)\\
Homework 5 (due Mar 13 Thu 10am)
\end{center}

\medskip\noindent
{\bf Instructions:} As in previous homeworks.



\begin{description}
\bigskip
\item[Problem 5.1:]
In this problem, we will explore other simple families of hash functions
and see how close they are to being universal.  Let $U\ge m$.

\begin{enumerate}
\item[(a)] (35 pts)
Let $p$ be a random prime in $[m/2,m]$.
Let $b$ be a random number in $\{0,\ldots,p-1\}$.
Define $h_p : \{0,\ldots,U-1\} \rightarrow \{0,\ldots,m-1\}$ by
\[ h_p(x) = x\bmod p.\]

Prove that for any fixed $x,y\in\{0,\ldots,U-1\}$ with $x\neq y$,
we have $\Pr_p[h_p(x)=h_p(y)] \le O((\log U)/ m)$.

(Thus, this hash function family is ``almost'' universal with an extra logarithmic factor.)

[\emph{Hint}: By the prime number theorem, the number of primes in $[m/2,m]$ is $\Theta(m/\log m)$.
Given a number $z\le U$, how many prime factors in $[m/2,m]$ can $z$ have?]

\medskip
\item[(b)] (35 pts)
Suppose $U=2^w$ and $m=2^\ell$.  We can view numbers in $\{0,\ldots,U-1\}$ as
binary vectors in $\{0,1\}^w$, and numbers in $\{0,\ldots,m-1\}$ as binary
vectors in $\{0,1\}^\ell$.

Pick a random 0-1 matrix $A$ of dimension $\ell \times w$ (each of its $\ell w$ entries is
an independently chosen random bit).
Define $h_A : \{0,1\}^w \rightarrow \{0,1\}^\ell$ by
\[ h_A(x) = (Ax) \bmod 2. \]
All arithmetic operations are done modulo 2; for example,
addition is equivalent to exclusive-or.

Prove that for any fixed $x,y\in\{0,1\}^w$ with $x\neq y$,
we have $\Pr_A[h_A(x)=h_A(y)] = 1/m$.

(Thus, this hash function family is universal, and its computation avoids integer division and requires only bitwise exclusive-or, which is cheaper.  However, the evaluation time isn't $O(1)$.)

[\emph{Hint}: suppose $x$ and $y$ differs at the $i$-th bit.  Focus on choices over the $i$-th column of $A$\ldots]

\medskip
\item[(c)] (30 pts)
Continuing part (b), let $b$ be a random 0-1 vector of dimension $\ell$.
Define $h_{A,b} : \{0,1\}^w \rightarrow \{0,1\}^\ell$ by
\[ h_{A,b}(x) = (Ax+b) \bmod 2. \]

Prove that for any fixed $x,y\in\{0,1\}^w$ with $x\neq y$ and for
any fixed $s,t\in \{0,1\}^\ell$,
we have $\Pr_{A,b}[(h_{A,b}(x)=s) \wedge (h_{A,b}(y)=t)] = 1/m^2$.

(Thus, this hash function family is 2-universal.) 

[\emph{Hint}: use (b).]
\end{enumerate}

\newpage
\textbf{Solution:}
\begin{enumerate}
  \item[(a)] 
    Let $ z = |x - y| $ for distinct $ x, y $. 
    A collision occurs if and only if $ p $ divides $ z $, i.e., 
    \begin{equation}
      h_p(x) = h_p(y) \quad \Leftrightarrow \quad p \mid z.
    \end{equation}
    
    The number of primes in $ [m/2, m] $ is given by the Prime Number Theorem as 
    \begin{equation}
      \Theta(m / \log m).
    \end{equation}
    
    If $ z $ has $ k $ prime factors in $ [m/2, m] $, then 
    \begin{equation}
      z \geq (m/2)^k.
    \end{equation}
    Since $ z \leq U $, we obtain
    \begin{equation}
      k \leq \log_{(m/2)} U.
    \end{equation}
    Thus, $ z $ has at most $ O(\log U) $ prime factors in $ [m/2, m] $.

    Since $ U \ge m $, the probability of choosing a prime $ p $ 
    that divides $ z $ is at most
    \begin{equation}
      \frac{O(\log U)}{\Theta(m / \log m)} = O\left(\frac{\log U}{m}\right).
    \end{equation}
  
  \item[(b)] 
    Let $ z = x - y $, the collision condition is:
    \begin{equation}
      h_A(x) = h_A(y) \quad \Leftrightarrow \quad A(x - y) = 0 \quad \Leftrightarrow \quad A z = 0.
    \end{equation}
    Since $ x \neq y $, we have $ z \neq 0 $, 
    meaning there exists at least one bit $ i $ 
    such that $ z_i = 1 $.

    Let $ a_1, a_2, \dots, a_w $ be the columns of $ A $, and express:
    \begin{equation}
      A z = z_1 a_1 + z_2 a_2 + \cdots + z_w a_w.
    \end{equation}
    Since $ z_i = 1 $ for at least one $ i $, rewrite:
    \begin{equation}
      A z = a_i + \sum_{j \neq i} z_j a_j.
    \end{equation}
    The elements of $ a_1, a_2, \dots, a_w $ are chosen independently 
    and uniformly at random from $ \{0,1\} $,
    thus each element of $\sum_{j \neq i} z_j a_j$ is independent and uniform.
    So the probability that $ A z = 0 $ is simply 
    the probability that each element of $ a_i = 0 - \sum_{j \neq i} z_j a_j$,
    this occurs with probability $ 1/2^\ell $.

    Thus, the probability of a collision is:
    \begin{equation}
      \Pr[A z = 0] = \frac{1}{2^\ell} = \frac{1}{m}.
    \end{equation}

  \item[(c)]
  Let $z = x-y$,
  \begin{equation}
    (h_{A,b}(x)=s) \wedge (h_{A,b}(y)=t)  \Rightarrow
    A(x-y) = s-t \Leftrightarrow
    Az = s-t
  \end{equation}
  For fixed $s$ and $t$, the probability that $ Az = s-t $ is 
  the probability that each element of $ a_i = (s-t) - \sum_{j \neq i} z_j a_j$, 
  this occurs with probability $1/2^\ell$:
  \begin{equation}
    \Pr_A[Az=s-t] = \frac{1}{m}
    \label{eq:5_1_c_1}
  \end{equation}
  Given $A$ s.t. $ Az = s-t $, 
  \begin{equation}
    (h_{A,b}(x)=s) \wedge (h_{A,b}(y)=t)  \Leftrightarrow
    b = Ax-s
  \end{equation}
  Since $b$ is a random 0-1 vector, we have
  \begin{equation}
    \Pr_b[b=Ax-s|Az=s-t] = \frac{1}{m}
    \label{eq:5_1_c_2}
  \end{equation}
  Combining \ref{eq:5_1_c_1} and \ref{eq:5_1_c_2}, we obtain:
  \begin{align}
    &\Pr_{A,b}[(h_{A,b}(x)=s) \wedge (h_{A,b}(y)=t)] \nonumber \\
    &= \Pr_A[Az=s-t] \Pr_b[b=Ax-s|Az=s-t] \nonumber \\
    &= \frac{1}{m^2}
  \end{align}
\end{enumerate}


\newpage
\item[Problem 5.2:]
We are given a subset $S$ of $\{1,\ldots,N\}$ and a number $k\le N$.
We want to compute a new set $f_k(S)=\{ a_1+\cdots+a_k: \mbox{$a_1,\ldots,a_k$ are $k$
\emph{distinct} elements of $S$}\}$.

(Note that this is a little different from one of the midterm 1 practice problems, which
is about computing the set $S+\cdots+S$ with $k$ terms, where an element may be used
more than once in the sum.)

\begin{enumerate}
\item[(a)] (50 pts)
First consider a related problem: given $\ell$ disjoint sets $S_1,\ldots,S_\ell\subseteq\{1,\ldots,N\}$
and a number $k\le\ell\le N$,
compute a new set $g_k(S_1,\ldots,S_\ell)=\{ a_1+\cdots+a_k: \mbox{$a_1,\ldots,a_k$ are elements chosen from $k$ \emph{different} sets among $S_1,\ldots,S_\ell$}\}$.

Give a (deterministic) algorithm that solves this problem in $O(\ell^3 N\log N)$ time or better.

[Hint: use divide-and-conquer, like in the midterm 1 practice question.
Recall that for $A,B\subseteq\{1,\ldots,M\}$, we can compute
$A+B=\{a+b: a\in A,\ b\in B\}$ in $O(M\log M)$ time by convolution/FFT\@.
It may be helpful to extend the problem to computing $g_k(S_1,\ldots,S_\ell)$ for all $k\le \ell$ simultaneously.]

\medskip
\item[(b)] (50 pts)
Now design and analyze a Monte Carlo algorithm to compute $f_k(S)$.  The running time should be of the form
$O(k^c N\log^{c'}N)$ for some constants $c,c'$, and the overall 
error probability should be at most $1/N$.

[Hint: partition $S$ into $\ell$ disjoint sets randomly or by using a hash function, and then use part (a).  How should we set $\ell$?]
\end{enumerate}


\newpage
\textbf{Solution:}
\begin{enumerate}
  \item[(a)]
  Split sets into two halves $L$ and $R$, compute all $g_k (0 \le k \le \ell)$ using divide-and-conquer. 
  For all possible $g_i(L)$ and $g_j(R)$ s.t. $i+j=k$, compute $g_i(L) + g_j(R)$ using FFT. 
  Take the union of all such sets into $g_k(S_1,\ldots,S_\ell)$.

  \begin{algorithm}
  \caption{}
  \begin{algorithmic}[1]
    \Function{DisjointSubsetSum} S
    \State $\ell = \text{length of }S$
    \If{$\ell=0$}
      \State \Return $\phi$
    \EndIf
    \State $L = S_1, S_2, \ldots, S_{\lfloor\ell/2\rfloor}$
    \State $R = S_{\lfloor\ell/2\rfloor+1}, \ldots, S_\ell$ 
    \State $G_L = \text{DisjointSubsetSum}(L)$
    \State $G_R = \text{DisjointSubsetSum}(R)$
    \State Initialize array $G$ of length $\ell + 1$, each element is $\phi$
    \For{$i$ = 0 to $\lfloor\ell/2\rfloor$}
      \For{$j$ = 0 to $\lfloor\ell/2\rfloor+1$}
        \If{$i+j\le\ell$}
          \State Compute $g = g_i(L) + g_j(R)$ by taking convolution of $G_L[i]$ and $G_R[j]$
          \State $G[i+j] = g \cup G[i+j]$
        \EndIf
      \EndFor
    \EndFor
    \State \Return $G$
    \EndFunction
    \State $S = [S_1, S_2, \ldots, S_\ell]$
    \State \Return DisjointSubsetSum(S)[k]
  \end{algorithmic}
  \end{algorithm}

  Time complexity analysis:

  Steps 1 and 2 split the problem into 2 subproblems and solve them recursively. 
  In step 3, both halves have $O(\ell)$ sets, 
  so there are $O(\ell^2)$ pairs of $(i,j)$ such that $i+j=k$.
  Since the largest number in any set is $\ell N$ and $\ell \le N$, 
  each set contains at most $O(\ell N)$ elements, 
  computing each $g_i(L) + g_j(R)$ using FFT takes $O(\ell N\log N)$ time. 
  Thus the total time complexity is:
  \begin{equation}
    T(\ell, N) = 2T(\ell/2, N) + O(\ell^3N\log N) = O(\ell^3N\log N).
  \end{equation}

  \item[(b)]
  Randomly partition $S$ in to $\ell$ disjoint sets.
  For any specific set of $k$ elements $(a_1, \ldots, a_k)$, the sum of the set $s$ appears
  in $g_k(S)$ if each element is properly assigned to different disjoint subset.
  Since $s$ can be computed by multiple sets, the possibility that $f_k(S)$ contains
  $s$ is:
  \begin{equation}
    q = \Pr[s\in g_k(S)] \ge \prod_{i=0}^{k-1} \frac{\ell-i}{\ell}.
  \end{equation}
  Let $\ell = k$, $q = \frac{k!}{k^k}$ is a constant for a particular $k$.
  Assume we randomly partition $S$ for $c$ times and compute $g_{k,i}(S)$ for the $i$th partition, 
  the possibility that $s$ does not appear in any $g_{k,i}(S)$ is:
  \begin{equation}
    \Pr\left[\bigcap_{i=1}^{c}s\notin g_{k,i}(S)\right] = (1-q)^c
  \end{equation}
  There are at most $kN$ elements in $f_k(S)$, accroding to Union Bound, 
  the possibility that the union of all $g_{k,i}(S)$
  does not contain all possible elements is:
  \begin{equation}
    \Pr\left[\bigcup_{s\in f_x(S)}\bigcap_{i=1}^{c}s\notin g_{k,i}(S)\right] 
    \le kN\cdot\Pr\left[\bigcap_{i=1}^{c}s\notin g_{k,i}(S)\right]
    = kN(1-q)^c
  \end{equation}
  Rewrite $kN(1-p)^c \le 1/N$, we obtain,
  \begin{equation}
    c \ge \frac{2logN}{|log(1-q)|}
  \end{equation}

  \begin{algorithm}
  \caption{}
  \begin{algorithmic}[1]
    \State $c$ = $\lceil2logN\rceil$
    \For{$i$ = 1 to $c$}
      \State Randomly partition $S$ into $k$ subsets
      \State Compute $g_{k,i}(S)$ using algorithm in (a)
    \EndFor
    \State \Return $\bigcup_{i=1}^{c}g_{k,i}(S)$
  \end{algorithmic}
  \end{algorithm}

  Let $c = O(logN)$, overall error probability is at most $1/N$, 
  total time complexity is $O(k^3Nlog^2N)$.



\end{enumerate}


\newpage
\item[Problem 5.3:]
Consider the following geometric problem: given a set $P$ of 
$n$ points in 2D, with integer coordinates from $\{0,1,\ldots,U-1\}$, 
find a {\em closest pair\/}, i.e., two points
$p,q\in P\ (p\neq q)$  with the smallest Euclidean distance.
Let $\delta(P)$ denote the distance of the closest pair.

We have seen an $O(n\log n)$-time divide-and-conquer algorithm from class.
In this question, we give a different, faster randomized algorithm
(which has the added advantage that it can be
extended to higher dimensions).

\begin{enumerate}
\item[(a)] (35 pts) First give an $O(n)$-expected-time (Las Vegas) algorithm for
 the easier \emph{decision problem}: given
a value $r$, decide whether $\delta(P)<r$.  

(Hints: Build a uniform grid where each cell is an
$(r/2)\times(r/2)$ square.  Use hashing.
How many points can a grid cell have?
For each grid cell, how many grid cells are of distance at most 
$r$?)

\medskip
\item[(b)] (65 pts) Now, consider the following recursive Las Vegas algorithm 
to compute $\delta(P)$:

\newcommand{\ClosestPair}{\textsc{Closest-Pair}}
\begin{quote}
\begin{tabbing}
19.\ \= for\= for\= for\= for\= for\=\kill
$\ClosestPair(P)$:\\[2pt]
1.\> if $|P|\le 100$ then return answer by brute force\\
2.\> partition $P$ into subsets $P_1,\ldots,P_{20}$ each with 
at most $\ceil{n/20}$ points\\
3.\> let $S=\{(i,j)\mid 1\le i<j\le 20\}$\\
4.\> $r=\infty$\\
5.\> for each $(i,j)\in S$ {\em in random order\/} do\\
6.\>\> if $\delta(P_i\cup P_j) < r$ then\\
7.\>\>\> $r=\ClosestPair(P_i\cup P_j)$\\
8.\> return $r$
\end{tabbing}
\end{quote}

Explain why the algorithm is always correct,
and analyze its expected running time by solving a recurrence.

(Hints: Where is part (a) used?  What is $|S|$?  What is the expected number of times 
line~7 is performed, using facts from class?)
\end{enumerate}


\newpage
\textbf{Solution:}
\begin{enumerate}
  \item[(a)] 
  Use a uniform grid to efficiently check for point pairs within distance $r$. 
  Divide the plane into a uniform grid where each cell is a square of size $(r/2) \times (r/2)$.
  If any cell contains two or more points, their distance is necessarily less than $r$.
  If a point has no neighbor in its own cell, 
  it only needs to check points in adjacent cells (at most 24 cells since the side length
  of the square is $r/2$).
  Use a hash table to store the occupied grid cells, ensuring $O(1)$ insertion and lookup in expectation.

  \begin{algorithm}
  \caption{}
  \begin{algorithmic}[1]
    \State Initialize a hash table $T$
    \For{each point $p = (x, y) \in P$}
      \State $(v,w) =
        \left(\left\lfloor \frac{x}{r/2} \right\rfloor, 
        \left\lfloor \frac{y}{r/2} \right\rfloor \right)$
      \State Transform and combine $v$ and $w$ into a 0-1 vector: $\{0,1\}^{2log_2U}$, 
      as the index of $p$.
      \State Store $p$ in the hash table $T$:
      \If{another point already exists in the cell}
        \State \Return True
      \Else
        \State Check all 24 neighboring cells (a $5\times5$ grid centered at $p$). 
        If any existing point in these cells is within distance $r$, \Return True.
      \EndIf
    \EndFor
    \State \Return False
  \end{algorithmic}
  \end{algorithm}

  Time complexity analysis:

  For each point, computing the grid index takes $O(1)$; 
  Hash table operations (insertion and lookup) take expected $O(1)$;
  Checking adjacent cells involves at most a constant number of comparisons, also $O(1)$.
  Thus, for $n$ points, the expected total runtime is $O(n)$.

  \item[(b)]
  The algorithm correctly finds the closest pair distance $\delta(P)$ 
  by employing a divide-and-conquer strategy with a randomized approach.

  Base Case: If $|P| \leq 100$, the algorithm directly computes the closest pair using brute force, 
  which is guaranteed to be correct.

  Partitioning: The set $P$ is divided into 20 subsets $P_1, P_2, \dots, P_{20}$, 
  each containing at most $\lceil n/20 \rceil$ points.

  The algorithm considers all pairs of subsets from the set 
    \[
    S = \{(i, j) \mid 1 \leq i < j \leq 20\}.
    \]
  Since there are 20 subsets, the number of pairs is:
    \[
    |S| = \binom{20}{2} = 190.
    \]

  The algorithm iterates over the pairs $(i,j)$ in random order. 
  If $\delta(P_i \cup P_j) < r$, it updates $r$ with the result of a recursive call 
  to $\text{ClosestPair}(P_i \cup P_j)$.
  The global closest pair distance $\delta(P)$ must exist in at least one of these subset pairs. 
  Thus, as long as the correct pair $(i^*, j^*)$ is considered, 
  the algorithm will correctly compute $\delta(P)$.

  The decision algorithm from part (a) is used to check whether $\delta(P_i \cup P_j) < r$ 
  before making recursive calls.

  Since the algorithm always considers a valid pair that contains the closest points, 
  it will return the correct answer.

  Time complexity analysis:

  Let $T(n)$ denote the expected runtime of the algorithm.
  The partitioning step (line 2) takes $O(n)$ time.
  There are 190 pairs, and each pair is checked in constant time (except for potential recursive calls).

  For each pair $(i,j)$, the decision problem (from part (a)) runs in $O(n)$ expected time. Since the subset $P_i \cup P_j$ has at most $O(n/10)$ points, the decision step per pair runs in $O(n)$.
  However, not all 190 pairs will result in recursive calls.
  Since the pairs are traversed in a random order, the optimal pair $(i^*, j^*)$ appears with equal probability at some position within the first $k$ iterations. Once the correct minimum distance $\delta(P)$ is found, subsequent pairs will not trigger recursive calls. 
  Thus the expected number of times line 7 (recursive call) executed is a constant.

  Therefore, in expectation, the algorithm makes a recursive call on a subset of at most $O(n/10)$ points, a constant number of times.

  The expected runtime satisfies the recurrence:
  \[
  T(n) = O(n) + O(1) \cdot T(n/10).
  \]

  Expanding the recurrence:
  \[
  T(n) = O(n) + T(n/10).
  \]
  Applying the recurrence iteratively:
  \[
  T(n) = O(n) + O(n/10) + O(n/10^2) + O(n/10^3) + \dots.
  \]

  This forms a geometric series:
  \[
  O\left(n \sum_{k=0}^{\log_{10} n} \frac{1}{10^k} \right).
  \]

  Since the sum of the geometric series converges to a constant, we get:
  \[
  T(n) = O(n).
  \]
\end{enumerate}


\end{description}
\end{document}



